% !TEX recipe = pdflatex only
\documentclass[10pt]{letter}

\usepackage[margin=1in, top=0.75in, bottom=0.75in]{geometry}
\usepackage{hyperref}
\usepackage{xcolor}
\usepackage{enumitem}
\hypersetup{colorlinks=true, urlcolor=blue, linkcolor=blue}

\signature{Hameem M.\ Mahdi, B.S.C.S., M.S.E., Ph.D.\\
Mayo Clinic, Rochester, MN 55905, USA\\
\href{mailto:mahdi.hameem@mayo.edu}{mahdi.hameem@mayo.edu} ~|~ ORCID: \href{https://orcid.org/0009-0007-0005-3080}{0009-0007-0005-3080}}

\address{Hameem M.\ Mahdi\\
Mayo Clinic\\
200 First Street SW\\
Rochester, MN 55905, USA}

\date{April 15, 2026}
\setlength{\parskip}{0.4em}

\begin{document}

\begin{letter}{Editorial Office\\
\emph{Artificial Intelligence}\\
Elsevier}

\opening{Dear Editors,}

I am pleased to submit ``Perceive, Plan, Act, Self-Correct: An Architectural Framework for Goal-Directed Agentic AI Systems'' for consideration as a Regular Article in \emph{Artificial Intelligence}.

\textbf{Summary.}\quad
The paper proposes the \textsc{Perceive--Plan--Act--Self-Correct} (PPAS) framework, a four-phase canonical loop for goal-directed agentic AI grounded in classical agent theory and extended for LLM-era architectures. We decompose agentic systems into an 8-layer technology stack mapping 15+ frameworks to a common reference, validated by three empirical studies: a benchmark meta-analysis (12 agents across SWE-Bench Verified, WebArena, and GAIA), a design-pattern comparison, and an inter-agent protocol evaluation (MCP, A2A, AG-UI).

\textbf{Key contributions:}
\begin{enumerate}[nosep, leftmargin=1.5em]
    \item A formal four-phase agent loop unifying existing architectural patterns;
    \item An 8-layer technology stack mapping 15+ frameworks to a common reference model;
    \item Empirical evidence that reflective planning with human-in-the-loop approval achieves the highest task-completion rates (72.5\% SWE-Bench Verified) while reducing irreversible-action failures by 73\%;
    \item A failure-mode taxonomy with EU AI Act governance mapping.
\end{enumerate}

\textbf{Significance for AIJ.}\quad
Agentic AI is one of the fastest-growing areas of AI research, yet the field lacks a unified architectural vocabulary. This work provides exactly such a unifying framework---formally grounded in classical agent theory (BDI, SOAR, OODA) and validated empirically. Agent architectures, planning, and multi-agent coordination are core topics of \emph{Artificial Intelligence}, making this a natural fit.

\textbf{Preprint.}\quad A preprint is available on engrXiv: \url{https://doi.org/10.31224/6738} (March 2026).

\textbf{Data and code availability.}\quad All datasets, scripts, and notebooks are at: \url{https://github.com/HAMEEMM/ai_systems_landscape_2026/tree/main/publications/01-agentic-ai}

The manuscript is not under consideration elsewhere. I confirm compliance with the journal's submission guidelines. Thank you for considering this submission.

\closing{Sincerely,}

\end{letter}

\end{document}
